\documentclass[11pt,a4paper]{appletmanual}
\usepackage[english]{babel}
\manuallang{en}

\appletname{Profit, Costs and Supply}
\appletsub{The competitive firm's short-run problem, and the supply curve that
comes out of it. Two panels drawn to the same scale, because they are the same
curve.}
\appletdir{firm-supply}

\begin{document}
\manualtitlepage{Applet manual · Introductory Microeconomics · Universidad de los Andes\\
Translated from the GeoGebra applet “Beneficios y Costos”.\\
MIT licensed. It opens in a browser; there is nothing to install.}

\manualtoc
\thispagestyle{empty}
\clearpage
\setcounter{page}{1}

% ==========================================================================
\section{What this is}

A competitive firm in the short run has exactly one degree of freedom: how much
to produce. Capital is fixed, the price is handed to it, and all it can do is
pick the quantity $q$ that maximises profit
\[
  \pi(q) \;=\; p\,q \;-\; c(q).
\]

That is the problem. The firm's supply curve is the \emph{answer} to that
problem, solved once for every possible price. They are two different objects,
and in almost every textbook they appear in two different figures on two
different pages, joined by a sentence saying that marginal cost ``is'' the
supply curve, which the student has to take on trust.

This applet puts them side by side and to the same scale. The right panel is
the problem; the left one is its answer. Marginal cost is drawn in both, in the
same coordinates, and in the supply panel it is laid over the supply curve as a
red dashed line. Above the shutdown price the dashes fall exactly on top.
Nothing has to be taken on trust.

\figher[1.0]{panels}{The two panels at the opening settings. Left: the firm's
supply curve in the $(q, p)$ plane. Right: the problem, with marginal, average
and average variable cost. Both axes are the same in both panels, so marginal
cost occupies the same place in each.}

\begin{amkey}
Marginal cost and inverse supply are not similar curves: above the bottom of
average variable cost they are the same curve, point for point. The whole
applet is built so that this can be seen rather than asserted.
\end{amkey}

\subsection{Who it is for}

A first course in microeconomics, at the session where profit maximisation
turns into the supply curve. It assumes the student has already met total,
average and marginal cost, and knows what an optimum is. It assumes no calculus
beyond the idea of a slope.

%% Sections shared by every manual, English edition.
%% Pulled in with \input{common-en}. The only thing that differs between
%% applets is \appletfolder, which is also the folder name and the path on the site.

\section{Opening it}

There is nothing to install. The whole application is one \code{index.html}
file: it downloads no libraries, calls no server, and stores nothing outside
your browser.

\subsection{The three ways}

\begin{description}
\item[Double-click the file.] The quickest route. Find
  \code{\appletfolder/index.html} and open it. The browser shows it straight off the
  disk, with an address beginning \code{file://}. It works offline from the
  first moment.
\item[From the web.] If someone has published it, the address ends in
  \code{/\appletfolder/}. Once the page has loaded you can disconnect: it needs the
  network for nothing else.
\item[Served from a folder.] For a class, putting the folder on any static
  server is enough. With Python installed:
  \begin{quote}\code{python3 -m http.server 8000}\end{quote}
  then \code{http://localhost:8000/\appletfolder/} in the browser.
\end{description}

\begin{amnote}
Saving the page with \key{Ctrl}\,+\,\key{S} (or \key{Cmd}\,+\,\key{S} on a
Mac) leaves a copy of your own that still works offline and that can travel on a
memory stick or go out by email.
\end{amnote}

\subsection{Language and theme}

Two pairs of buttons sit at the top right.

\begin{ctrltable}{Control & What it does}
ES / EN & Switches the whole interface between Spanish and English on the spot.
          Nothing is recomputed, so you can switch mid-explanation without
          losing the state you had. \\
Dark / Light & Switches the theme. Dark is the original design; light exists
          for projectors and for printing. \\
\end{ctrltable}

Both choices are remembered in the browser, so next time it opens the way you
left it.

\subsection{Which browser}

Any reasonably recent one: Chrome, Edge, Firefox or Safari, on a computer,
tablet or phone. On a narrow screen the panels stack one above the other
instead of sitting side by side. No account, no permissions, no extensions.


% ==========================================================================
\section{The screen}

\fig[1.0]{interface}{The whole screen. Controls down the left; the verdict
along the top; the two panels with their readouts below.}

Left to right, top to bottom:

\begin{description}
\item[The control rail.] Picks the cost function, moves its parameters, sets
  the price, decides what the firm panel shows and which layers are drawn.
  Every slider has a numeric box beside it: drag, or type the exact number.
\item[The verdict.] The coloured band under the title. It names which case the
  firm is in --- profit, loss, shutdown, break-even, capacity, no optimum ---
  and explains why. It changes on its own as you move anything.
\item[The supply panel] (left). The inverse supply curve in the $(q, p)$ plane,
  with the shutdown price, the break-even price and the current price.
\item[The firm panel] (right). Either the marginal curves or the totals,
  depending on the \ui{Firm panel} switch.
\item[The readouts.] The row of numbers under each panel. They are the same
  values that are drawn, written to three decimals for when the eye is not
  enough.
\end{description}

\subsection{Dragging to change the price}

In any panel whose vertical axis \emph{is} the price --- the supply panel
always, the firm panel when it is showing marginals --- you can drag with the
mouse or a finger and the price follows the pointer. It is the natural way to
read a supply curve: pick a price, see what quantity comes out.

In the totals view the vertical axis is money, not money per unit, and the
panel does not answer the pointer.

\subsection{Keyboard}

\begin{ctrltable}{Key & What it does}
\key{1} \key{2} \key{3} \key{4} & Pick the cost function. \\
\key{m} & Firm panel to marginals. \\
\key{t} & Firm panel to totals. \\
\end{ctrltable}

\subsection{The dark theme}

Dark is the original design; light is there for projectors and for printing.
The two were designed separately --- it is not an inversion of the colours.

\fig[0.95]{dark}{The same two panels in the dark theme.}

% ==========================================================================
\section{The model}

\subsection{Costs}

Everything starts from a total cost function $c(q)$: what it costs to produce
$q$ units when the capital is already paid for. The rest follow from it.

\begin{align*}
  \text{Fixed cost}          &&  F   &= c(0) \\
  \text{Variable cost}       &&  VC(q) &= c(q) - F \\
  \text{Marginal cost}       &&  MC(q) &= c'(q) \\
  \text{Average cost}        &&  AC(q) &= \frac{c(q)}{q} \\
  \text{Average variable cost}&& AVC(q) &= \frac{c(q) - F}{q}
\end{align*}

Fixed cost is read as the limit of $c(q)$ as $q$ falls to zero from above. In
the families that start from a production function that limit is $r\bar K$, the
rental on the fixed capital: it is paid whatever is produced.

\subsection{The problem}

\[
  \max_{q \,\ge\, 0} \;\; \pi(q) \;=\; p\,q - c(q).
\]

If the maximum is interior, the first-order condition gives
$MC(q^\ast) = p$, and the second-order condition requires marginal cost to be
\emph{rising} there. Both matter: on a cubic cost curve the equation
$MC(q) = p$ has two roots and one of them is a \emph{minimum} of profit.

But $q^\ast = 0$ is a candidate too, and not a limiting case of the other one:
it is a corner, with profit $\pi(0) = -F$. The firm produces only if producing
beats closing.

\subsection{The shutdown price}

Comparing the two:
\[
  p\,q - c(q) \;>\; -F
  \quad\Longleftrightarrow\quad
  p\,q \;>\; c(q) - F = VC(q)
  \quad\Longleftrightarrow\quad
  p \;>\; AVC(q).
\]

Fixed cost drops out of the comparison. It is already sunk, and a decision that
takes it into account is a decision taken wrongly. Out of this comes the
\textbf{shutdown price}:
\[
  \underline{p} \;=\; \min_{q > 0} AVC(q),
\]
and the rule everybody recites: the firm produces as long as the price covers
average variable cost, even when it does not cover average total cost.

\subsection{The break-even price}

\[
  p_{be} \;=\; \min_{q>0} AC(q).
\]
The price at which profit is exactly zero. Above it there is profit; below it
there are losses --- which may still be worth taking, as long as the price
stays above $\underline{p}$.

\subsection{Supply}

\[
  q^\ast(p) \;=\;
  \begin{cases}
    0 & \text{if } p < \underline{p},\\[2pt]
    \text{the root of } MC(q) = p \text{ on the rising branch} & \text{if } p > \underline{p}.
  \end{cases}
\]

At $p = \underline{p}$ the firm is indifferent, and supply \textbf{jumps}: it
goes from zero to the quantity $\underline{q}$ that minimises average variable
cost, without passing through the values in between. That jump is not a
technicality. It is why the supply curve of a firm with fixed costs does not
reach the origin.

\subsection{Producer surplus}

\[
  PS \;=\; \int_0^{q^\ast}\!\bigl(p - MC(q)\bigr)\,dq
       \;=\; p\,q^\ast - VC(q^\ast)
       \;=\; \pi + F.
\]

One number written three ways, and the applet draws two of them: in the right
panel as the area between $p$ and marginal cost, and in the left as the area
between the price axis and the supply curve, from $\underline{p}$ up to $p$.
The two regions have different shapes --- with a jump they cannot coincide ---
and exactly the same area.

\begin{amnote}
Producer surplus and profit are the same thing only when there is no fixed
cost. With $F > 0$ the surplus is always larger, by exactly $F$.
\end{amnote}

% ==========================================================================
\section{The four cost functions}

Picked with the four buttons at the top of the rail, or with \key{1} to
\key{4}. Each brings its own sliders.

\subsection{1 · Quadratic / power}

\[
  c(q) = F + a\,q^{k}, \qquad k > 1.
\]
The textbook shape. Marginal cost, $MC = a k q^{k-1}$, rises from zero, so
average variable cost has no interior minimum: its infimum is zero and is
approached only in the limit. The consequence: \textbf{this firm never shuts
down}. At any positive price it is worth producing something.

\begin{ctrltable}{Slider & What it is}
$F$ & Fixed cost. Raises or lowers average cost without touching marginal
      cost. \\
$a$ & Scale of variable cost. \\
$k$ & Curvature. With $k$ near 1 marginal cost is nearly constant and the
      problem stops having a solution; with $k$ large it rises very fast and
      supply goes vertical. \\
\end{ctrltable}

\subsection{2 · Custom}

Type $c(q)$ into the text box. Marginal cost is obtained by differentiating the
expression \emph{symbolically}, not numerically, so it is exact. The one
supplied by default,
\[
  c(q) = 0.1\,q^{3} - 0.5\,q^{2} + 2q,
\]
is the textbook cubic: no fixed cost, U-shaped average variable cost with
minimum $1.375$ at $q = 2.5$, and therefore a strictly positive shutdown price
and a visible jump in supply.

\fig[1.0]{cubic}{The cubic at $p = 1.6$. The supply panel shows all three
pieces of the curve: flat against the price axis below $1.375$, the
discontinuous jump, and the rising branch that coincides with marginal cost.}

\subsection{3 · Cobb--Douglas production}

In the short run capital is fixed at $\bar K$ and the only variable factor is
labour. Starting from
\[
  f(K, l) = b\,\bigl(K^{\beta} l^{1-\beta}\bigr)^{s},
\]
invert it for the labour needed to produce $q$,
\[
  l(q) = \Bigl(\frac{q}{b\,\bar K^{\beta s}}\Bigr)^{\!e}, \qquad
  e = \frac{1}{(1-\beta)s},
\]
and price it at the wage $w$:
\[
  c(q) \;=\; r\bar K \;+\; w\,b^{-e}\,\bar K^{-\beta/(1-\beta)}\,q^{\,e}.
\]

Everything turns on $e$. With $e > 1$ there are diminishing returns to the
variable factor, marginal cost rises, and the problem has a solution. With
$e \le 1$ marginal cost does not rise, every extra unit adds profit, and there
is no optimal quantity. The applet says so instead of drawing something
meaningless.

\fig[0.8]{no-optimum}{The verdict when $e \le 1$: no optimum, and what to move
so that there is one.}

\begin{ctrltable}{Slider & What it is}
$b$ & Total factor productivity. \\
$\beta$ & Capital's elasticity. It enters $e$. \\
$s$ & Returns to scale of the Cobb--Douglas. It enters $e$. \\
$w$ & The wage. Scales all of variable cost. \\
$r$ & The rental on capital. It is the fixed cost, $F = r\bar K$. \\
$\bar K$ & The fixed capital. It is what makes the short run short. \\
\end{ctrltable}

\subsection{4 · Logistic log-CD production}

The same idea, with a production function that \emph{saturates}:
\[
  g(K, l) = \frac{b\,K^{\gamma}}{1 + e^{-s(\ln f - m)}},
\]
whose ceiling is $\bar Q = b\,\bar K^{\gamma}$ however much labour is hired.
Inverting it,
\[
  c(q) \;=\; r\bar K \;+\; w\,e^{m/(1-\beta)}\,\bar K^{-\beta/(1-\beta)}
             \left(\frac{q}{\bar Q - q}\right)^{\!e},
  \qquad 0 < q < \bar Q .
\]

Marginal cost blows up as $\bar Q$ is approached, and the supply curve has a
vertical asymptote there. It is the honest picture of what a fixed capital
means: however far the price rises, there is a quantity this plant cannot pass.

\fig[1.0]{capacity}{The logistic family at $p = 3.2$. The capacity
$\bar Q = b\bar K^{\gamma}$ is marked by a grey vertical, and marginal cost
climbs towards it without ever arriving.}

% ==========================================================================
\section{The controls}

\subsection{Price}

A slider from 0 to 10, with a numeric box for an exact value. Under it a
sentence says where the price stands relative to the shutdown and break-even
prices. It can also be moved by dragging inside any panel whose vertical axis
is the price.

\subsection{Firm panel}

\begin{ctrltable}{View & What it draws}
Marginals & $MC$, $AC$ and $AVC$, in money per unit. That is the same unit as
  the price, so this panel shares \emph{both} axes with the supply panel and
  the two figures are comparable pixel for pixel. \\
Totals & $c(q)$, revenue $pq$ and profit $\pi(q)$, in money. The vertical axis
  is no longer the supply panel's; the quantity axis still is. The pill in the
  panel header says which of the two cases you are in. \\
\end{ctrltable}

\subsection{The layers}

The list changes with the view, because not every layer makes sense in both.

\subsubsection{In marginals}

\begin{ctrltable}{Layer & What it draws}
Marginal cost & $MC(q)$, in red. It also appears laid over the supply panel. \\
Average cost & $AC(q)$, in gold. Its minimum is the break-even price. \\
Average variable cost & $AVC(q)$, in orange. Its minimum is the shutdown
  price. \\
Price line & The horizontal at $p$. \\
Variable-cost area & The region under $MC$ out to $q^\ast$, which is
  $VC(q^\ast)$. \\
Producer surplus area & The region between $p$ and $MC$ out to $q^\ast$. \\
Optimum & The point $(q^\ast, p)$, its vertical, and the profit rectangle
  between $p$ and $AC(q^\ast)$ --- green for profit, red for a loss. \\
Break-even point & Marks the minima of $AC$ and of $AVC$. \\
Grid & The background rule. \\
\end{ctrltable}

\fig[1.0]{areas}{Both areas switched on at once, with the cubic at $p = 1.6$.
On the left the surplus measured as the integral of $q$ with respect to price;
on the right the same number as the integral of $p - MC$ with respect to
quantity.}

\subsubsection{In totals}

\begin{ctrltable}{Layer & What it draws}
Total cost & $c(q)$, with a dotted horizontal at the fixed cost. \\
Total revenue & The line $pq$. \\
Profit & $\pi(q) = pq - c(q)$. Its maximum is at $q^\ast$. \\
Variable cost & $c(q) - F$, the same curve dropped. \\
Tangent of slope $p$ & The line tangent to $c(q)$ at $q^\ast$. It is parallel
  to the revenue line: there is the first-order condition, drawn. \\
Optimum & The vertical segment between $c(q^\ast)$ and $pq^\ast$, whose length
  \emph{is} the profit. \\
Break-even point & The quantities where $\pi(q) = 0$. \\
\end{ctrltable}

\fig[1.0]{totals}{The totals view with the cubic at $p = 1.6$. The dotted
tangent touches the cost curve exactly where its slope equals the price, and
the vertical segment at $q^\ast$ measures the profit.}

% ==========================================================================
\section{How the optimum is found}

Worth knowing, because the applet does not do what the GeoGebra original did
and the difference shows in the interesting cases.

\textbf{It does not solve $MC = p$.} That equation has two roots on a cubic
cost curve, and one of them is a minimum of profit. With any non-convex cost
function the global optimum can additionally sit at a corner, where no
first-order condition finds it.

What it does instead is maximise profit directly:

\begin{enumerate}
\item Bound the range from above, by finding the point past which marginal cost
  has overtaken the price and is still rising.
\item Scan that range with 700 points and keep the best.
\item Refine that neighbourhood by golden section.
\item Compare the result against shutting down, which gives $\pi(0) = -F$.
  Shutting down is a candidate in its own right, not a limiting case.
\end{enumerate}

The supply curve is traced the same way: a full maximisation at each of 140
prices. It is the only way the jump comes out. Inverting $MC = p$ would never
produce it, because below the shutdown price the answer is zero and at the
shutdown price it leaps.

\begin{amnote}
The minimum of average variable cost is searched over a geometric ladder, not a
linear one: what decides whether the firm ever shuts down is the shape of the
curve near zero, and a uniform mesh over five orders of magnitude resolves
nothing there. If the minimum lands on an end of the ladder then there is no
minimum --- the curve is still falling --- and the applet reports that there is
no shutdown price instead of quoting the smallest value it happened to sample.
\end{amnote}

% ==========================================================================
\section{The verdict}

The band at the top names the case and explains it. There are six.

\begin{ctrltable}{Verdict & When}
Produces at a profit & $p > p_{be}$. The price beats minimum average cost. \\
Produces at a loss & $\underline{p} < p < p_{be}$. And is right to: closing
  would cost the whole fixed cost. \\
Zero profit & $p = p_{be}$, the bottom of average cost. \\
Shuts down & $p < \underline{p}$. The optimum is $q^\ast = 0$ and the loss is
  exactly $F$. \\
Against the capacity & The optimum is practically at $\bar Q$. Only the
  logistic family reaches this. \\
No optimum & Marginal cost never overtakes the price over any range worth
  drawing. \\
\end{ctrltable}

\fig[1.0]{shutdown}{Shutdown, with the cubic at $p = 1.15$. The supply curve
lies flat against the price axis, $q^\ast$ is shown in red, and the verdict
explains why producing would be worse than not producing.}

% ==========================================================================
\section{Classroom activities}

Each one starts from the opening settings unless it says otherwise. If you get
lost, reload the page.

\begin{activity}{The two curves are one}
\amset Cost function \ui{Quadratic / power}, opening values, $p = 2$. Switch
\ui{Marginal cost} on if it is off.
\amexp The cyan curve in the left panel and the red one in the right occupy the
same place on the screen. Move $p$: the purple point travels along both at once
and always at the same height. Raise $a$ or $k$ and both steepen together.
\par\smallskip
Ask the class why the left panel has no average cost curve drawn in bold:
because average cost plays no part in the decision of \emph{how much} to
produce, only in whether the business is worth being in at all.
\end{activity}

\begin{activity}{Manufacturing a shutdown price}
\amset Choose \ui{Custom}. The default expression is the cubic.
\amexp $AVC$ is now U-shaped with minimum $1.375$, and supply no longer starts
at the origin: there is a stretch flat against the price axis and a jump.
\par\smallskip
Take the price down to $1.30$ and then up to $1.40$. The quantity goes from $0$
to $2.55$ in one step. Ask: what quantity does the firm produce when
$p = 1.375$ exactly? The honest answer is that it does not care, and that is
why supply there is not a function but a correspondence.
\end{activity}

\begin{activity}{Producing at a loss}
\amset Cubic. This cubic has $F = 0$, so its shutdown and break-even prices
coincide at $1.375$ and there is no loss band. Give it a fixed cost: type
\code{1+0.1q\^{}3-0.5q\^{}2+2q} into the box. Set $p = 1.6$.
\amexp The shutdown price is still $1.375$ --- fixed cost does not touch
average variable cost --- but break-even rises to $1.733$. There is now a band
of prices between them, and $1.6$ is inside it.
\par\smallskip
The \ui{Profit} readout shows $-0.394$ in red, and still the optimum is to
produce $2.869$. Ask why it does not close. Closing would cost the whole fixed
cost, $-1$; producing loses only $0.394$. Sweep the price across the whole band
$1.375 < p < 1.733$: profit is negative throughout and producing stays best.
\end{activity}

\begin{activity}{Surplus measured twice}
\amset Cubic, $p = 1.6$. Switch on \ui{Producer surplus area}.
\amexp Two shaded regions of clearly different shapes appear, one in each
panel. The \ui{Surplus} readout gives $0.606$, and it is both of theirs.
\par\smallskip
Check it with the numbers in the right panel: $q^\ast = 2.869$,
$\text{revenue} = p\,q^\ast = 4.59$, $c(q^\ast) = 3.98$. Since $F = 0$ here,
variable cost is the whole cost, and $PS = 4.59 - 3.98 = 0.61$. It matches both
areas.
\par\smallskip
Ask why the two regions are not the same shape if they measure the same thing.
Because one integrates $p - MC$ with respect to quantity and the other
integrates $q^\ast$ with respect to price; with a jump at the shutdown price
they cannot be the same region, but the fundamental theorem of calculus says
they have the same area.
\end{activity}

\begin{activity}{The first-order condition, drawn}
\amset Cubic, $p = 1.6$, firm panel on \ui{Totals}.
\amexp The dotted tangent to total cost at $q^\ast$ is parallel to the revenue
line. Move the price: the revenue line pivots, and the point of tangency slides
along the cost curve to wherever the slope matches again.
\par\smallskip
The vertical segment at $q^\ast$ is the profit. Switch on \ui{Break-even point}
and the two quantities appear at which that segment has zero length.
\end{activity}

\begin{activity}{Two roots, one optimum}
\amset Cubic, $p = 1.6$, \ui{Marginals} view.
\amexp The price line cuts marginal cost \emph{twice}. The purple point is on
the right-hand crossing, where $MC$ is rising.
\par\smallskip
Ask what happens at the other one. It is the solution of $MC = p$ that a
student would get by solving the equation without checking the second
derivative, and it is a minimum of profit: producing that quantity is the worst
the firm can do among those that cover variable cost. Switch to \ui{Totals} to
see it as a valley in the profit curve.
\end{activity}

\begin{activity}{Capacity}
\amset \ui{Logistic log-CD production} family, opening values.
\amexp The grey vertical marks $\bar Q = b\bar K^{\gamma}$. Push the price to
the top of its slider: the quantity grows by less and less, and the
\ui{Elasticity} readout approaches zero.
\par\smallskip
Now raise $\bar K$ and the capacity moves right. That is exactly the long-run
exercise the short run forbids: the only way to produce past $\bar Q$ is to
change plant.
\end{activity}

\begin{activity}{When there is no optimum}
\amset \ui{Cobb--Douglas production} family. Take $\beta$ down to $0.05$ and
$s$ up to $2$.
\amexp The formula in the rail shows $e = 0.526$, and the verdict changes to
``No optimum''. Every readout in the supply panel turns into a dash.
\par\smallskip
Marginal cost is falling: each extra unit costs less than the last and sells
for the same price. Ask what would happen to such a firm in a competitive
market. The answer --- that it would grow until it stopped being competitive
--- is the argument for why diminishing returns are not a technical assumption
but the condition for the model to mean anything.
\end{activity}

% ==========================================================================
\section{Expression syntax}

In the cost box the variable is \code{q}. \code{x} is accepted too, because
that is what GeoGebra called it.

\begin{ctrltable}{Element & What is allowed}
Operators & \code{+ - * / \^{} ( )} \\
Functions & \code{sqrt ln log log10 exp abs sin cos tan min max pow} \\
Constants & \code{e}, \code{pi} \\
Implicit multiplication & \code{2q}, \code{3(q+1)}, \code{2sqrt(q)} \\
Associativity & \code{\^{}} binds to the right and tighter than unary minus:
  \code{-q\^{}2} is $-(q^2)$ \\
\end{ctrltable}

Marginal cost is differentiated symbolically. Where the expression has a kink
--- \code{min}, \code{max}, \code{abs} --- there is no symbolic derivative to
be had and it falls back to a central difference, which is what the picture
shows anyway.

\begin{amwatch}
Typing a decreasing cost function, or one with constant marginal cost, lands
you in the ``No optimum'' case. That is not a failure of the applet: that
problem has no solution.
\end{amwatch}

% ==========================================================================
\section{Changes from the GeoGebra original}

\begin{description}
\item[Two panels instead of one.] The original drew total cost, revenue,
  profit, marginal cost and average cost on one pair of axes. Money and money
  per unit share no scale, so the marginal curves were unreadable next to the
  totals. They are separated here, and the marginal panel is locked to the
  supply panel's scale.
\item[The supply curve, which was not there.] Along with the shutdown price,
  the break-even price, the jump between them, and the elasticity.
\item[The optimum is a maximisation, not an intersection.] The original
  intersected \code{Eq: CM(x) = p} with the axis and used the result as the
  right end of a \code{Max[Beneficios, 0, x(A)]}. With non-monotone marginal
  cost that bracket is the wrong one, and when $MC = p$ has no root the whole
  chain returned \code{NaN} and the applet went blank.
\item[Shutting down is a candidate.] The original always compared against an
  interior maximum. Here $\pi(0) = -F$ competes with it, which is what makes
  the shutdown price exist.
\item[An infimum is not a minimum.] Average variable cost that merely decreases
  towards zero has no minimum, so there is no shutdown price. Reporting the
  smallest sampled value gave a shutdown price of $7.7\times10^{-4}$.
\item[Sliders bounded away from division by zero.] $\beta = 1$ divides by
  $1-\beta$; $\bar K = 0$ raises zero to a negative power; $w = 0$ and $a = 0$
  make cost constant; $k = 1$ makes marginal cost constant. All were reachable
  in the original.
\item[$m$ widened] from $[0, 1]$ to $[-1, 2]$, where it actually moves the cost
  level enough to see.
\item[Bilingual, and with a verdict.] Every state of the problem is named and
  explained, because the shapes on their own do not say which of the six cases
  a student is looking at.
\item[Dead objects not carried over.] \code{AreaProfi}, \code{CM2} (a copy of
  \code{CM}), \code{eq1}, \code{A}, \code{prueba} and \code{prueba2} were
  plumbing for the GeoGebra construction rather than objects with meaning.
\end{description}

% ==========================================================================
\section{Troubleshooting}

\begin{ctrltable}{Symptom & What is happening}
The panels are empty & The cost expression does not parse. The red message
  under the box says whether it is a syntax problem or a variable problem. \\
Everything reads ``$\infty$'' or a dash & You are in the ``No optimum'' case.
  Raise $k$, or $s$, or $\beta$, or switch family. \\
The scales stop matching & The firm panel is on \ui{Totals}. The pill in its
  header says so. Go back to \ui{Marginals}. \\
The axis moves while I drag the price & Deliberate: axes are rounded to round
  numbers and step in jumps, so that two settings can be compared by eye
  instead of renumbering continuously. \\
The price ignores the mouse & You are dragging inside the totals panel, whose
  vertical axis is money. Use the supply panel or the slider. \\
The language or theme was lost & They are kept in the browser. In a private
  window, or after clearing site data, they go back to their defaults. \\
\end{ctrltable}

%% Sections shared by every manual, English edition.
%% Pulled in with \input{common-en}. The only thing that differs between
%% applets is \appletfolder, which is also the folder name and the path on the site.

\section{Publishing it for a class}

Any static host will do: GitHub Pages, Netlify, a university web directory,
even a shared folder served over HTTP. The only thing to upload is the
\code{index.html} file, inside a folder named however you want the address to
read.

The repository ships a GitHub Actions workflow
(\code{.github/workflows/pages.yml}) that publishes every applet at once on each
push to the default branch. To switch it on, in the repository:
\emph{Settings\,$\rightarrow$\,Pages\,$\rightarrow$\,Source: GitHub Actions}.
That is the one step the workflow cannot do for itself, because creating a Pages
site needs administration rights that a \code{GITHUB\_TOKEN} is never granted.

\begin{amnote}
To hand the applet to students without reliable internet, send them the
\code{index.html} file directly. It weighs less than a photograph and opens with
a double-click.
\end{amnote}

\section{Licence and provenance}

The application code is MIT. Use it, change it and pass it on, in class or out
of it, with attribution.

This applet is a standalone rewrite of an original GeoGebra applet. The rewrite
is not a literal translation: where the original had a construction error, a
division by zero reachable with its own sliders, or a dead object inherited from
whatever applet it was copied from, this version fixes it and says so. The
section \emph{Changes from the GeoGebra original} lists those changes one by
one.

Everything is hand-written and dependency-free: the expression parser, the
symbolic differentiation, the level-curve tracing and the drawing. In
particular there is no \code{eval}, so what a student types into a text box
cannot become code, and a strict content security policy does not break the
page.


\end{document}
